\documentclass[11pt]{article}

\usepackage[a4paper,margin=1in]{geometry}
\usepackage{amsmath,amssymb,amsthm,mathtools}
\usepackage[hidelinks]{hyperref}
\hypersetup{
  pdfauthor={Bangyen Pham},
  pdftitle={Displaced-Zero Certificates and the Coefficient Mass of Polynomial Multiples},
  pdfsubject={Lower bounds for the coefficient order statistics and logarithmic coefficient mass of polynomial multiples},
  pdfkeywords={polynomial multiples, coefficient mass, reduced height, Chebyshev systems, exponential sums},
  pdflang={en-US}
}

\newtheorem{theorem}{Theorem}[section]
\newtheorem{lemma}[theorem]{Lemma}
\newtheorem{corollary}[theorem]{Corollary}
\newtheorem{proposition}[theorem]{Proposition}
\newcommand{\R}{\mathbb R}
\newcommand{\N}{\mathbb N}
\DeclareMathOperator{\Tail}{Tail}
\DeclareMathOperator{\sgn}{sgn}

\title{Displaced-Zero Certificates and the Coefficient Mass of Polynomial Multiples}
\author{\shortstack{Bangyen Pham\\
\small Independent Researcher\\
\small\href{mailto:bangyenp@gmail.com}{\texttt{bangyenp@gmail.com}}\\
\small\href{https://orcid.org/0009-0006-9928-2088}{ORCID: 0009-0006-9928-2088}}}
\date{September 2026}

\begin{document}
\maketitle

\begin{abstract}
Fix distinct real roots $2\le r_1<\cdots<r_L$.  For every real polynomial
$F=\sum_jf_jx^j$ divisible by $\prod_{i=1}^L(x-r_i)$ with $|f_D|\ge1$ we prove
a degree-independent bound on each coefficient order statistic: the
$(u+1)$-st largest nonleading coefficient magnitude is at least
$\prod_{i=u+1}^L(r_i-1)$.  Summing bounds the logarithmic coefficient mass
$\Lambda(F)=\sum_j\log^+|f_j|$; for the first $L$ primes this gives
$\Lambda(F)\ge(\tfrac12+o(1))L^2\log L$, with asymptotic equality for
$F=\prod_{i=1}^L(x-p_i)$.  The engine is a
prefix-sensitive comparison of the infinite $\ell^1$ tails of normalized
exponential sums with prescribed positive-integer zeros, strict unless the
zeros are consecutive.  We isolate the zero count behind the comparison as an
interlacing lemma.  The argument builds on the reduced-height certificate
method of Dubickas and Jankauskas.
\end{abstract}

\noindent\textbf{Keywords.} Polynomial multiples; coefficient mass; reduced
height; Chebyshev systems; exponential sums.

\noindent\textbf{2020 Mathematics Subject Classification.} Primary 11C08.

\section{Introduction}

For $F(x)=\sum_j f_jx^j$, define its logarithmic coefficient mass by
\[
  \Lambda(F)=\sum_{j:f_j\ne0}\log^+|f_j|,
  \qquad \log^+x=\max(0,\log x),
\]
where all logarithms are natural.  Throughout, $\N_{>0}=\{1,2,\ldots\}$, and
every asymptotic statement is for $L\to\infty$ with absolute implied
constants unless parameters are named.

Reduced length and reduced height minimize the two classical coefficient
norms over monic multiples~\cite{mignotte,dubickas,schinzel,schinzel-complex,dj},
the Euclidean analogue is effectively computable~\cite{frw}, and the dual
problems---a multiple of bounded height, or a root of prescribed multiplicity
for a restricted-coefficient power series---form a classical
line~\cite{bombieri-vaaler,bbbp-multiple,bbbp-ray}.  A bound on the height
$\max_j|f_j|$ concerns a single coefficient, and the length $\sum_j|f_j|$ can
be carried by many comparable terms, so neither norm alone forces a hierarchy
of large coefficient magnitudes.  The mass $\Lambda$
instead accumulates over all coefficients of magnitude exceeding one: it is
the logarithm of the multiplicative height,
$\Lambda(F)=\log\prod_{j:f_j\ne0}\max(1,|f_j|)$.  In a fixed base
$\Lambda(F)/\log B$ is the total base-$B$ digit length of the nonzero
coefficient digits, so it agrees with the encoded length of the coefficient
vector up to a term proportional to the number of encoded coefficients; for
real coefficients it is a mass rather than a literal encoding length.
It is additive over coefficients, and this is what makes it the natural
target here.  A multiplicative bound on one coefficient order statistic
contributes one additive term to $\Lambda$; a bound holding for successive
order statistics therefore accumulates into a lower bound for $\Lambda$ that
no statement about a single height parameter can produce.  The results below
are in this sense distributional: many coefficients must be large, not merely
the largest one.

Dubickas and Jankauskas associate to reciprocal roots and a prescribed zero
set $J$ a normalized determinant sequence $S_n(J)$ with $S_0(J)=1$ and
$S_k(J)=0$ for $k\in J$~\cite{dj}.  That sequence is the exponential sum $u$
studied below: its infinite tail is the reciprocal of their dual lower bound
for the reduced height, their inequality~(8), and they evaluate the
consecutive set $J=\{1,\ldots,c-1\}$
explicitly~\cite[Section~2]{dj}.  Our new ingredient is a displaced-zero
comparison for the $\ell^1$-tail of this certificate: moving a prescribed
integer zero and the largest reciprocal node together strictly decreases the
tail, and fixed-node displacement alone need not.  To our knowledge, earlier
work uses these certificates to control a single height parameter, whereas the
iterative exclusion argument below produces simultaneous lower bounds for
successive coefficient order statistics.  The comparison is strict, with
equality exactly for the consecutive zero set.  Sparse-multiple work
asks for multiples under prescribed degree, sparsity and height
constraints~\cite{grt}; our bound is uniform in the degree and in all
additional roots of the multiple.

The main result is a distributional bound on the coefficient sizes of a
multiple.  For distinct positive roots $2\le r_1<\cdots<r_L$,
Theorem~\ref{thm:order}, via the certificate Lemma~\ref{lem:slack}, bounds all
$L$ coefficient order statistics of every multiple, and
Corollary~\ref{cor:mass} sums them into a lower bound for $\Lambda$:
\[
  b_k(F)\ge\prod_{i=k}^{L}(r_i-1),
  \qquad
  \Lambda(F)\ge\sum_{i=1}^{L}i\log(r_i-1),
  \qquad 1\le k\le L .
\]
Its engine is Theorem~\ref{thm:tail}, the comparison behind everything: it
bounds the tail of a certificate by the product over its leading consecutive
zeros and shows the bound is attained only there.  Corollary~\ref{cor:prime}
specializes to prime roots, Corollary~\ref{cor:infimum} sharpens the resulting
asymptotic to the infimum over normalized multiples, and
Proposition~\ref{prop:confluent} removes the distinctness of the nodes.

\section{A displaced-zero tail theorem}

Let $\rho_1>\cdots>\rho_c\ge2$ and $y_i=\rho_i^{-1}$, so that
$0<y_1<y_2<\cdots<y_c\le1/2$.  For any $Z\subset\N_{>0}$ of size $c-1$, the
generalized Vandermonde system
\[
  u_d=\sum_{i=1}^c a_i y_i^d,
  \qquad u_0=1,
  \qquad u_z=0\quad(z\in Z),
\]
determines the coefficients $a_i$ uniquely.  Indeed, a kernel vector would be
a nonzero $c$-node exponential sum with the $c$ zeros $\{0\}\cup Z$,
contradicting the zero bound below.  It is convenient to regard the sequence
$u_0,u_1,\ldots$ as the restriction to $d\in\{0,1,2,\ldots\}$ of the
real-analytic function $d\mapsto\sum_ia_iy_i^d$; all statements about zeros
and signs of an exponential sum refer to this function of the real variable.

\paragraph{Zero bound.}
A nonzero exponential sum on $r$ distinct positive nodes has at most $r-1$
real zeros, counted with multiplicity.  This is classical: divide by the
power of the largest node so the sum becomes a real combination of
$e^{\lambda d}$ with distinct $\lambda\le0$; the $r$ exponentials span a
Chebyshev system, so by Rolle's theorem applied inductively to successive
derivatives the sum has at most $r-1$ zeros, the final derivative being a
nonzero multiple of $e^{\lambda_1d}$ and hence zero-free~\cite[Chapter~I,
Theorem~4.1]{karlin-studden}; it is also the exponential form of Laguerre's
extension of Descartes' rule of signs~\cite[Part~V]{polya-szego}.
Consequently an $r$-node
sum with $r-1$ prescribed real zeros has exactly those zeros, each is simple,
and the sum changes sign at each of them.

\begin{theorem}[Displaced-zero comparison]\label{thm:tail}
Let
$f=\max\{0\le k\le c-1:\{1,\ldots,k\}\subseteq Z\}$.  Then
\begin{equation}\label{eq:tail}
  \Tail(u):=\sum_{d\ge1}|u_d|
  \le \prod_{i=1}^{f+1}\frac{y_i}{1-y_i}
  =\frac{1}{\prod_{i=1}^{f+1}(\rho_i-1)}.
\end{equation}
Equality holds exactly when $Z=\{1,\ldots,c-1\}$.
\end{theorem}

For a concrete two-node example, take $(\rho_1,\rho_2)=(3,2)$.  If the
prescribed zero is consecutive, $Z=\{1\}$, then
\[
 u_d=3^{1-d}-2^{1-d},\qquad \Tail(u)=\frac12,
\]
so equality holds in \eqref{eq:tail}.  Moving the zero to $Z=\{2\}$ gives
\[
 u_d=\frac95\,3^{-d}-\frac45\,2^{-d},\qquad \Tail(u)=\frac3{10}<\frac12.
\]
Here displacement lowers the tail, but that is not the general behaviour: at
fixed nodes the comparison is not monotone.  For $(\rho_1,\rho_2)=(5,3)$,
$Z=\{1\}$ gives $\Tail(u)=1/8$ while $Z=\{2\}$ gives $\Tail(u)=9/64>1/8$, so
the displaced tail can exceed the consecutive one.  What the theorem proves is
not fixed-node monotonicity but a comparison in which the largest prescribed
zero and the largest reciprocal node are removed together, reducing the node
count.  Each displaced zero creates an extra sign change relative to a
comparison sum, and the proof shows that this change reduces the total tail.
It treats the displaced zeros one at a time, peeling the last one.

The hard step compares a correction sum at the marked points of the zero set
with its value just above them.  We state it separately; the proof is the
zero-count argument that was implicit in the comparison.

\begin{lemma}[Interlacing of the correction sum]\label{lem:interlace}
Let $y_1<\cdots<y_c\le1/2$, let $Z'\subset\N_{>0}$ satisfy $|Z'|=c-2$, put
$m=\max(Z'\cup\{0\})$, and fix $z>m$.  Let $F$ be the sum on
$y_1,\ldots,y_{c-1}$ with $F_0=1$ and $F|_{Z'}=0$, and $G$ the sum on
$y_1,\ldots,y_c$ with $G_z=1$ and $G|_{\{0\}\cup Z'}=0$.  Write
$\widehat F_d=\sgn(F_z)F_d$, $E_d=G_d-y_cG_{d-1}$,
$\sigma=G_{m+1}/\widehat F_{m+1}$, and
\[
  \Delta_d=\sigma\widehat F_d-E_d.
\]
Then $\Delta_d\ge0$ for every $d\ge m+1$.
\end{lemma}

\begin{proof}
The sums $E$ and $\widehat F$ both live on the $c-1$ nodes
$y_1,\ldots,y_{c-1}$: expanding the definition of $G$ gives
\[
  E_d=\sum_{i=1}^{c-1}a_i(y_i-y_c)y_i^{\,d-1}
      =\sum_{i=1}^{c-1}b_iy_i^d,
\]
with $b_i=a_i(1-y_c/y_i)$, the $i=c$ term having cancelled.  Hence $\Delta$
is an exponential sum on $c-1$ nodes, so by the zero bound it is either
identically zero or has at most $c-2$ real zeros counted with multiplicity.

Relabel $A=\{0\}\cup Z'$ as $a_1=0<a_2<\cdots<a_{c-1}=m$; let $r$ be the
number of consecutive pairs $\{a_j,a_{j+1}\}=\{d-1,d\}$ in $A$, and
$J=\{j:a_{j+1}\in Z',\ a_j<a_{j+1}-1\}$, so that $|J|=c-2-r$.

By the zero bound, $G$ has exactly the simple zeros $\{0\}\cup Z'$; on a
component $(a_j,a_{j+1})$ its sign is $(-1)^{c-1-j}$ because it is positive
past the last zero and flips at each.  We record the zeros of $\Delta$.

First, $\Delta_{m+1}=0$, because $G_m=0$ and $\sigma=G_{m+1}/\widehat
F_{m+1}$.  Second, if $a_j=a_{j+1}-1$ then $a_j,a_{j+1}\in A$, so
$G_{a_j}=G_{a_{j+1}}=0$, and $a_{j+1}\in Z'$, so $\widehat F_{a_{j+1}}=0$;
thus $\Delta_{a_{j+1}}=-E_{a_{j+1}}=-G_{a_{j+1}}+y_cG_{a_{j+1}-1}=0$.  There
are $r$ such forced zeros, one for each consecutive pair.  At every other
marked point $a_{j+1}\in Z'$ one has $a_j<a_{j+1}-1$, so
$a_{j+1}-1\notin A$ and $G_{a_{j+1}-1}\ne0$, and
\[
  \Delta_{a_{j+1}}=-E_{a_{j+1}}=y_cG_{a_{j+1}-1},
  \qquad \sgn\Delta_{a_{j+1}}=(-1)^{c-1-j}.
\]

If $J=\varnothing$ then $r=c-2$, and the forced zeros together with
$m+1$ give $r+1=c-1>c-2$ zeros, so $\Delta\equiv0$ and the claim is
immediate.  Assume $J\ne\varnothing$ and order $J=\{j_1<\cdots<j_p\}$, so
$p=|J|=c-2-r$.  Let $N_\Delta$ be the number of zeros of $\Delta$ counted
with multiplicity.  Take consecutive $j<j'$ in $J$.  Strictly between
$a_{j+1}$ and $a_{j'+1}$ lie $q=j'-j-1$ marked points of $A$, all with
indices outside $J$, hence $q$ forced zeros; and the endpoint signs
$(-1)^{c-1-j}$, $(-1)^{c-1-j'}$ differ by $(-1)^{q+1}$.  A real function
whose endpoint values differ in sign by $(-1)^{q+1}$ has an odd or even
number of zeros in the open interval according as $q+1$ is odd or even, so
the number of zeros in $(a_{j+1},a_{j'+1})$, counted with multiplicity, has
the parity of $q+1$ and is at least $q$; hence it is at least $q+1$, one more
than the $q$ forced zeros.  The $r$ forced zeros lie at marked points, and the
parity supplies one additional unit of multiplicity in each of the $p-1$ gaps
between consecutive members of $J$, possibly at a point already forced.  With
the zero at $m+1$ this gives
\[
  N_\Delta\ge r+(p-1)+1=c-2,
\]
which exhausts the zero-multiplicity budget of the $c-1$-node sum $\Delta$; in
particular there are no other zeros.

Every one of the $p-1$ extra multiplicities lies below $a_{j_p+1}$, so every
zero of $\Delta$ above $a_{j_p+1}$ belongs to the exhaustion just described
and is simple.  There are $c-1-j_p$ of them: the forced zeros at
$a_{j+1}$ for $j_p<j\le c-2$, and the zero at $m+1$.  Since
$\sgn\Delta_{a_{j_p+1}}=(-1)^{c-1-j_p}$ and there is no zero above
$a_{j_p+1}$ but those $c-1-j_p$ simple ones, crossing them flips the sign to
$(-1)^{2(c-1-j_p)}=+1$, and it stays positive.  Thus $\Delta_d\ge0$ for
$d\ge m+1$.
\end{proof}

\begin{proof}[Proof of Theorem~\ref{thm:tail}]
Induct on the number $c-1-f$ of displaced zeros.

\emph{Consecutive base case.}  For $Z=\{1,\ldots,c-1\}$, put
\[
  U(x)=\sum_{d\ge0}u_dx^d=\frac{N(x)}{\prod_i(1-y_ix)},
  \qquad \deg N\le c-1.
\]
The first $c-1$ nonconstant coefficients vanish, so $N(x)-\prod_i(1-y_ix)$
is a constant multiple of $x^c$; that constant is fixed by the vanishing
$x^c$ coefficient of $N$ (which has degree at most $c-1$), giving
\[
  N(x)=\prod_i(1-y_ix)+(-1)^{c+1}\Bigl(\prod_i y_i\Bigr)x^c .
\]
Since $1/\prod_i(1-y_ix)=\sum_{j\ge0}h_j(y_1,\ldots,y_c)x^j$, it follows that
for $d\ge c$,
\[
  u_d=(-1)^{c+1}\Bigl(\prod_i y_i\Bigr)h_{d-c}(y_1,\ldots,y_c),
\]
where $h_j$ is the complete homogeneous symmetric polynomial.  These terms
are one-signed, so
\[
  \Tail(u)=\Bigl(\prod_i y_i\Bigr)\sum_{j\ge0}h_j(y_1,\ldots,y_c)
          =\prod_i\frac{y_i}{1-y_i},
\]
which is \eqref{eq:tail} with $f=c-1$.

\emph{Displaced step.}  Suppose $f<c-1$, set $z=\max Z$ (displaced, hence
$z\ge f+2$), $Z'=Z\setminus\{z\}$, and $m=\max Z'$, taking $m=0$ if $Z'$ is
empty.  Let $F$ be the sum on $y_1,\ldots,y_{c-1}$ with $F_0=1$ and
$F|_{Z'}=0$, and $G$ the sum on all $c$ nodes with
$G|_{\{0\}\cup Z'}=0$ and $G_z=1$.  Then
\[
  u=F+\lambda G,\qquad \lambda=-F_z,
\]
because both sides satisfy the same $c$ interpolation conditions.  The set
$Z'$ has the same leading run $f$ as $Z$, and $f+1\le c-1$, so the induction
hypothesis bounds $\Tail(F)$ by the right side of \eqref{eq:tail}.  It
remains to prove $\Tail(u)<\Tail(F)$.

\emph{Sign comparison.}  By the zero bound, $G$ has exactly the simple zeros
$\{0\}\cup Z'$ and $F$ exactly $Z'$; below $z$ the sum $u$ has the same simple
zeros $Z'$ and no others, and at $z$ it has one additional simple zero.  Write
$\eta=\sgn(F_z)$ and $\mu=|F_z|$, so $u=F-\eta\mu G$.  Since $F$ and $G$ flip
at exactly the same points of $Z'$, their relative sign
$\sgn(F_d/G_d)=\eta$ is constant off $Z'$; thus $F_d=\eta\sgn(G_d)|F_d|$ and
$G_d=\sgn(G_d)|G_d|$, whence
\[
  u_d=\eta\sgn(G_d)\bigl(|F_d|-\mu|G_d|\bigr),
  \qquad
  |u_d|=\bigl||F_d|-\mu|G_d|\bigr|.
\]
The sign of $|F_d|-\mu|G_d|$ is therefore that of $\eta\sgn(G_d)u_d$.  Passing
a point of $Z'$ flips both $\sgn u$ and $\sgn G$, leaving this sign unchanged,
so it is constant on $(0,z)\setminus Z'$; it is positive at $d=0$, where
$u_0=F_0=1$ and $G_0=0$.  At $z$, where $G$ does not vanish, only $\sgn u$
flips, so the sign is negative for all $d>z$.  Hence
\[
  |F_d|>\mu|G_d|\quad(1\le d<z,\ d\notin Z'),
  \qquad
  |F_d|<\mu|G_d|\quad(d>z),
\]
with equality exactly at the zeros of $u$: at the points of $Z'$ both sides
vanish, and at $z$ one has $|F_z|=\mu|G_z|$.  For an exponential sum $V$, write
\[
  \Gamma_V(z)=\frac{\sum_{d\ge z}|V_d|}{|V_z|},
  \qquad A_G=\sum_{1\le d<z}|G_d|.
\]
Splitting $\Tail(F)-\Tail(u)$ at $z$ and using the two displayed
inequalities, $|F_z|=\mu|G_z|$ and $|G_z|=1$, gives
\begin{equation}\label{eq:split}
  \Tail(F)-\Tail(u)
  =\mu\bigl(2\Gamma_F(z)-\Gamma_G(z)+A_G\bigr).
\end{equation}

\emph{The correction ratio.}  We claim
\begin{equation}\label{eq:gamma}
  \Gamma_G(z)\le\frac{\Gamma_F(z)}{1-y_c}\le2\Gamma_F(z).
\end{equation}
Let $\widehat F_d=\eta F_d>0$ for $d>m$ and
$E_d=G_d-y_cG_{d-1}$, an exponential sum on the first $c-1$ nodes.  We first
show
\begin{equation}\label{eq:drive}
  \frac{E_d}{\widehat F_d}\le
  \frac{G_z}{\widehat F_z}\qquad(d\ge z>m).
\end{equation}
Set $A=\{0\}\cup Z'=\{a_1=0<\cdots<a_{c-1}=m\}$,
$\sigma=G_{m+1}/\widehat F_{m+1}$, and
$\Delta_d=\sigma\widehat F_d-E_d$.  Lemma~\ref{lem:interlace} applies to this
data and gives $\Delta_d\ge0$ for $d\ge m+1$, that is,
$E_d/\widehat F_d\le\sigma$ for $d\ge m+1$.  Moreover the $c$-node sum
\[
  W_d=\widehat F_{m+1}G_d-G_{m+1}\widehat F_d
\]
vanishes on $Z'\cup\{m+1\}$, which is $c-1$ points, so by the zero bound
those are its only zeros; beyond $m+1$ it has constant sign, and that sign is
positive because asymptotically only $G$ contains the largest node $y_c$,
with positive coefficient.  Hence $G_d/\widehat F_d>\sigma$ for $d>m+1$,
while $\sigma=G_z/\widehat F_z$ when $z=m+1$.  Therefore
\[
  \frac{E_d}{\widehat F_d}\le\sigma\le\frac{G_z}{\widehat F_z}
  \qquad(d\ge z>m),
\]
which is \eqref{eq:drive}.

Induction in $G_d=y_cG_{d-1}+E_d$, with $G_z=1$ and
\eqref{eq:drive}, gives for $t\ge0$
\[
  \frac{G_{z+t}}{G_z}
  \le\sum_{j=0}^t y_c^{t-j}
       \frac{\widehat F_{z+j}}{\widehat F_z}.
\]
Summing over $t$ (a geometric series in $y_c$) proves the first inequality
in \eqref{eq:gamma}; the second uses $y_c\le1/2$.

\emph{Strictness.}  It remains to see that $A_G>0$, so that
\eqref{eq:split} is positive.  By the zero bound, $G$ has exactly the simple
zeros $\{0\}\cup Z'$ and none other.  If $z>m+1$, then $m+1<z$ and
$m+1\notin\{0\}\cup Z'$, so $G_{m+1}\ne0$ and $A_G>0$.  If $z=m+1$, then the
displacement means $Z'$ does not contain all of $\{1,\ldots,z-1\}$, so some
$d$ with $1\le d<z$ avoids $\{0\}\cup Z'$, and again $A_G>0$.  Thus the
difference in \eqref{eq:split} is strictly positive, $\Tail(u)<\Tail(F)$,
and the induction closes.  Equality in \eqref{eq:tail} can occur only in the
base case.
\end{proof}

For example, when $c=2$, $\rho_1=3$, $\rho_2=2$, and $Z=\{1\}$, solving
$u_0=1,u_1=0$ gives $u_d=3(1/3)^d-2(1/2)^d$.  Hence
$\Tail(u)=1/((3-1)(2-1))=1/2$, illustrating both the consecutive-zero
formula and its equality case.

The restriction $\rho_i\ge2$ is substantive.  For $c=2$,
$y_1=1/10$, $y_2=4/5$, and $Z=\{2\}$, the normalized coefficients are
$a_1=64/63$ and $a_2=-1/63$, so
$\Tail(u)=364/2835>1/9=y_1/(1-y_1)$; Proposition~\ref{prop:cutoff} shows that
no cutoff larger than $1/2$ is universal.

\section{Coefficient order statistics}

Let
\[
 P(x)=\prod_{i=1}^L(x-r_i),\qquad 2\le r_1<\cdots<r_L,
\]
and let $F(x)=\sum_{j=0}^Df_jx^j\in\R[x]$ be a nonzero multiple of $P$.  For a
nonzero coefficient vector write $b_1\ge b_2\ge\cdots$ for the magnitudes of
the coefficients $f_j$ with $j<D$, in decreasing order.

\begin{lemma}[Slack certificate]\label{lem:slack}
Assume $|f_D|=1$.  Fix $0\le u\le L-1$.  For every set $U$ of $u$
coefficient positions below $D$, some $j<D$ outside $U$ satisfies
\[
 |f_j|\ge\prod_{i=u+1}^L(r_i-1).
\]
\end{lemma}

\begin{proof}
Use all $c=L$ roots $\rho_1=r_L>\rho_2=r_{L-1}>\cdots>\rho_c=r_1$.  A
position $j$ has backward distance $d=D-j$.  Let $W\subseteq\N_{>0}$ be the
set of distances of the positions in $U$, and let $w$ be the certificate on
these $c$ nodes, with $w_0=1$ and with zero set $Z$ consisting of $W$
together with the first $L-u-1$ positive distances outside $W$.  The $u$
positions have distinct distances, so $|W|=u$, and $L-u-1\ge0$ because
$u\le L-1$; hence $|Z|=c-1$, as required; write $w_d=\sum_i a_iy_i^d$.

The dual-certificate identity is explicit.  Since $F(\rho_i)=0$ for the $c$
chosen roots,
\[
  0=\sum_{i=1}^c a_i\rho_i^{-D}F(\rho_i)
   =\sum_{j\le D}f_jw_{D-j}
   =f_D+\sum_{j<D}f_jw_{D-j},
\]
using $w_0=1$.  The terms with $j\in U$ vanish because $D-j\in W\subseteq Z$.
Hence, as $|f_D|=1$,
\begin{equation}\label{eq:dual}
  1\le\sum_{j<D,\ j\notin U}|f_j|\,|w_{D-j}|
  \le\Bigl(\max_{j<D,\ j\notin U}|f_j|\Bigr)
        \sum_{j<D,\ j\notin U}|w_{D-j}|.
\end{equation}

The set $Z$ contains $\{1,\ldots,L-u-1\}$: a distance $t\le L-u-1$ is in
$W\subseteq Z$ if occupied by $U$, and otherwise is among the first
$L-u-1$ distances outside $W$.  Thus the leading run of $Z$ is at least
$L-u-1$, and Theorem~\ref{thm:tail} bounds the tail by the product over the
$L-u$ largest nodes.  These are $\rho_1,\ldots,\rho_{L-u}$, that is,
$r_L,r_{L-1},\ldots,r_{u+1}$; the last index is
$L-(L-u)+1=u+1$, which is why the product below starts there.  Therefore
\[
  \sum_{j<D,\ j\notin U}|w_{D-j}|
  \le\sum_{d\ge1}|w_d|=\Tail(w)
  \le\prod_{i=1}^{L-u}\frac{\rho_i^{-1}}{1-\rho_i^{-1}}
  =\frac{1}{\prod_{i=u+1}^L(r_i-1)}.
\]
Substituting into \eqref{eq:dual} gives
$1\le\bigl(\max_{j<D,j\notin U}|f_j|\bigr)/\prod_{i=u+1}^L(r_i-1)$, which is
the claim.
\end{proof}

\begin{theorem}[Coefficient order statistics]\label{thm:order}
If $|f_D|\ge1$, then for every $0\le u\le L-1$,
\begin{equation}\label{eq:order}
  b_{u+1}\ge\prod_{i=u+1}^L(r_i-1).
\end{equation}
\end{theorem}

\begin{proof}
Assume first $|f_D|=1$.  The assertion that some $j<D$ outside an arbitrary
set $U$ of $u$ positions has $|f_j|\ge T$ is equivalent to the assertion that
at least $u+1$ positions $j<D$ have $|f_j|\ge T$, that is, to
$b_{u+1}\ge T$: any $U$ of size $u$ misses one of the $u+1$ large positions,
and conversely the $u$ largest positions form a set $U$ that the
$(u+1)$-st largest avoids.  Apply this with $T=\prod_{i=u+1}^L(r_i-1)$ to
the conclusion of Lemma~\ref{lem:slack}.

For $|f_D|\ge1$, divide by $f_D$.  The quotient has leading coefficient $1$
and coefficient magnitudes $b_k(F)/|f_D|$, so
$b_k(F)\ge b_k(F/f_D)\ge T$ for every $k$.
\end{proof}

\begin{corollary}[Logarithmic mass]\label{cor:mass}
If $|f_D|\ge1$, then
\begin{equation}\label{eq:mass}
 \Lambda(F)\ge
 \sum_{u=0}^{L-1}
 \sum_{i=u+1}^L\log(r_i-1).
\end{equation}
\end{corollary}

\begin{proof}
Each product in \eqref{eq:order} is at least one because $r_i\ge2$, so the
$(u+1)$ largest magnitudes are all at least one.  Summing the logarithms of
\eqref{eq:order} over $u$ and observing that
$\Lambda(F)\ge\sum_{u=0}^{L-1}\log b_{u+1}$, the terms on
the right being a subset of $\sum_j\log^+|f_j|$ in decreasing order,
gives \eqref{eq:mass}.
\end{proof}

\begin{corollary}[Prime roots]\label{cor:prime}
Assume $|f_D|\ge1$ and that the increasing roots $r_i$ are bounded below by
the $i$th prime.  Then
\[
 \Lambda(F)\ge\Bigl(\frac12+o(1)\Bigr)L^2\log L .
\]
In particular $\Lambda(F)=\Omega(L^2\log L)$, and the bound is asymptotically
attained: for $F=P_L=\prod_{i=1}^L(x-p_i)$,
$\Lambda(P_L)\le(\tfrac12+o(1))L^2\log L$.
\end{corollary}

\begin{proof}
Write $p_i$ for the $i$th prime, so $r_i\ge p_i\sim i\log i$ and, since
$r_i-1\ge p_i-1$,
$\log(r_i-1)\ge\log p_i+O(1/p_i)$.  A term with index $i$ occurs in
\eqref{eq:mass} once for each $u$ with $u+1\le i$, that is, exactly $i$ times.
Hence
\[
 \Lambda(F)\ge\sum_{i=1}^L i\log(r_i-1)
 \ge\sum_{i=1}^L i\bigl(\log i+\log\log i+o(1)\bigr),
\]
and $\sum_{i\le L}i\log i=(\tfrac12+o(1))L^2\log L$, so the remaining terms
are $o(L^2\log L)$.  This is $\bigl(\tfrac12+o(1)\bigr)L^2\log L$.

For the matching upper bound, the nonleading coefficients of $P_L$ are the
elementary symmetric sums $e_k(p_1,\ldots,p_L)\le\binom{L}{k}p_L^k$, so
\[
 \Lambda(P_L)\le\sum_{k=1}^L\log\binom{L}{k}
 +\frac{L(L+1)}2\log p_L
 \le L^2\log2+\frac{L(L+1)}2\log p_L,
\]
and $\frac{L(L+1)}2\log p_L=(\tfrac12+o(1))L^2\log L$ while
$L^2\log2=o(L^2\log L)$.
\end{proof}

\begin{corollary}[Sharp prime-root infimum]\label{cor:infimum}
Let $M_L=\inf\Lambda(F)$, the infimum taken over the normalized multiples of
$P_L=\prod_{i=1}^L(x-p_i)$.  Then
\[
  M_L=\frac12L^2\log L+\frac12L^2\log\log L+O(L^2).
\]
\end{corollary}

\begin{proof}
The lower bound of Corollary~\ref{cor:prime} keeps the $\log\log i$ term:
\[
  \Lambda(F)\ge\sum_{i\le L}i\log(p_i-1)
  =\frac12L^2\log L+\frac12L^2\log\log L+O(L^2),
\]
since $\sum_{i\le L}i\log i=\frac12L^2\log L+O(L^2)$ and
$\sum_{i\le L}i\log\log i=\frac12L^2\log\log L+o(L^2)$.  The matching upper
bound is the one for $P_L$ from Corollary~\ref{cor:prime},
\[
  \Lambda(P_L)\le L^2\log2+\frac{L(L+1)}2\log p_L
  =\frac12L^2\log L+\frac12L^2\log\log L+O(L^2).
\]
\end{proof}

\begin{proposition}[Confluent analogue]\label{prop:confluent}
Let $0<y_1<\cdots<y_c\le1/2$ and $e_1,\ldots,e_c\ge1$, put $m=e_1+\cdots+e_c$,
and let $y_{(1)}\le\cdots\le y_{(m)}$ be the multiset with each $y_i$ repeated
$e_i$ times.  Let
\[
  u_d=\sum_{i=1}^c P_i(d)y_i^d,\qquad \deg P_i<e_i,
\]
satisfy $u_0=1$ and $u_z=0$ for a set $Z$ of $m-1$ positive integers with
leading run $f$.  Then
\[
  \Tail(u)\le\prod_{j=1}^{f+1}\frac{y_{(j)}}{1-y_{(j)}}.
\]
\end{proposition}

\begin{proof}
The interpolation system is indexed by the Hermite basis, not by the
ordinary exponentials, and the two must be kept apart.  For
$\varepsilon>0$, $i=1,\ldots,c$ and $k=0,\ldots,e_i-1$ put
\[
  \phi_{i,k,\varepsilon}(d)
  =y_i^d\Bigl(\frac{1-e^{-\varepsilon d}}{\varepsilon}\Bigr)^k .
\]
For fixed $\varepsilon>0$, the binomial expansion gives
\[
  \phi_{i,k,\varepsilon}(d)
  =\frac1{\varepsilon^k}\sum_{j=0}^{k}(-1)^j\binom{k}{j}
    \bigl(y_ie^{-j\varepsilon}\bigr)^d,
\]
so each $\phi_{i,k,\varepsilon}$ is a finite linear combination of the
ordinary exponentials $d\mapsto z^d$ at the distinct nodes
$z_{i,j}=y_ie^{-j\varepsilon}$.  For $0<\varepsilon<\varepsilon_0$ the $m$
nodes $z_{i,j}$ are distinct, lie in $(0,1/2]$, and stay separated between
distinct $i$.  Within one group $i$ the change of basis from
$\{z_{i,j}^{d}\}_{j<e_i}$ to
$\{\phi_{i,k,\varepsilon}\}_{k<e_i}$ is triangular with diagonal entries
$(-1)^k\varepsilon^{-k}$, hence invertible for $\varepsilon>0$, and the same
function space is spanned either way.

For such an $\varepsilon$ let $u^{(\varepsilon)}$ be the unique exponential
sum on the nodes $z_{i,j}$ with $u^{(\varepsilon)}_0=1$ and
$u^{(\varepsilon)}_z=0$ for $z\in Z$; existence and uniqueness are the
Vandermonde argument opening Section~2.  Theorem~\ref{thm:tail} applies with
the same $f$ and gives
\[
  \Tail(u^{(\varepsilon)})
  \le\prod_{j=1}^{f+1}\frac{y_{(j)}(\varepsilon)}{1-y_{(j)}(\varepsilon)},
\]
where $y_{(1)}(\varepsilon)\le\cdots\le y_{(m)}(\varepsilon)$ is the sorted
perturbed multiset.

It remains to pass to the limit.  Write
$a^{(\varepsilon)}=(a^{(\varepsilon)}_{i,k})$ for the coefficients of
$u^{(\varepsilon)}$ in the $\phi$-basis; they solve
$M(\varepsilon)a^{(\varepsilon)}=e_0$, where $M(\varepsilon)$ is the
$m\times m$ matrix with rows indexed by $d\in\{0\}\cup Z$ and columns by
$(i,k)$, with entry $\phi_{i,k,\varepsilon}(d)$, and $e_0$ is the first
standard basis vector.  Each entry
$\phi_{i,k,\varepsilon}(d)=y_i^d\bigl((1-e^{-\varepsilon d})/\varepsilon\bigr)^k$
has a removable singularity at $\varepsilon=0$ and extends analytically there
with limit $d^ky_i^d$, so $M(\varepsilon)\to M(0)$ as $\varepsilon\to0$, where
$M(0)$ is the confluent matrix.  The functions
$d\mapsto d^ky_i^d$ form an extended Chebyshev system on $\R$, for which
confluent Hermite interpolation is unisolvent~\cite[Chapter~I, Theorem~4.1]
{karlin-studden}; equivalently $M(0)$ is invertible.  Matrix inversion is
continuous on the open set of invertible matrices, so
$a^{(\varepsilon)}\to M(0)^{-1}e_0$, and for every fixed $d$,
\[
  u^{(\varepsilon)}_d
  =\sum_{i,k}a^{(\varepsilon)}_{i,k}\phi_{i,k,\varepsilon}(d)
  \longrightarrow\sum_{i,k}a_{i,k}d^ky_i^d=u_d .
\]
Fatou's lemma for series, applied to the nonnegative terms
$|u^{(\varepsilon)}_d|$, gives
\[
  \Tail(u)=\sum_{d\ge1}|u_d|
  \le\liminf_{\varepsilon\to0}\Tail(u^{(\varepsilon)}).
\]
Finally the order statistics of a converging multiset converge, so the
product $\prod_{j=1}^{f+1}y_{(j)}(\varepsilon)/(1-y_{(j)}(\varepsilon))$
tends to $\prod_{j=1}^{f+1}y_{(j)}/(1-y_{(j)})$.
\end{proof}

\begin{proposition}[The node cutoff $1/2$ is sharp]\label{prop:cutoff}
No constant $a>1/2$ can replace $1/2$ in Theorem~\ref{thm:tail} or in
Proposition~\ref{prop:confluent}: for every $a\in(1/2,1)$ some configuration
with all nodes in $(0,a]$ violates the bound
$\Tail(u)\le\prod_{j=1}^{f+1}y_{(j)}/(1-y_{(j)})$.
\end{proposition}

\begin{proof}
Take a single node $a\in(1/2,1)$ of multiplicity $c$ and
$Z=\{2,3,\ldots,c\}$, so the sum is $u_d=a^d\prod_{k=2}^c(1-d/k)$.  Then
$u_0=1$, the prescribed zeros are exactly $Z$, the leading run is $f=0$, and
\[
  u_1=\frac ac,\qquad
  |u_d|=\frac1c\binom{d-2}{c-1}a^d\quad(d\ge c+1),
\]
so, summing the binomial series,
\[
  \Tail(u)=\frac ac+\frac{a^{c+1}}{c(1-a)^c}
  =\frac ac\Bigl[1+\Bigl(\frac a{1-a}\Bigr)^c\Bigr].
\]
The bound with $f=0$ would be $a/(1-a)$, but
\[
  \frac{\Tail(u)}{a/(1-a)}
  =\frac{1-a}{c}\Bigl[1+\Bigl(\frac a{1-a}\Bigr)^c\Bigr]\longrightarrow\infty
\]
because $a/(1-a)>1$, so $\Tail(u)>a/(1-a)$ for all large $c$.

To transfer the obstruction to distinct nodes, suppose the bound held for
every configuration with nodes in $(0,a]$.  Coalescing the $c$ nodes near $a$
as in the proof of Proposition~\ref{prop:confluent} gives sums
$u^{(\varepsilon)}$ on distinct nodes in $(0,a]$ that satisfy it, and
$u^{(\varepsilon)}_d\to u_d$ for each $d$.  By Fatou's lemma,
$\Tail(u)\le\liminf_{\varepsilon\to0}\Tail(u^{(\varepsilon)})\le a/(1-a)$,
contradicting the previous paragraph for large $c$.
\end{proof}

\section{Scope}

Corollary~\ref{cor:mass} concerns normalized coefficient magnitudes, not
ordinary length or height.  Subject to $|f_D|\ge1$, it is uniform over the
degree of $F$ and all additional real or complex roots, and bounds all $L$
order statistics rather than about half of them.  Corollary~\ref{cor:prime}
shows the resulting prime-root bound $(1/2+o(1))L^2\log L$ is attained by
$P_L=\prod_{i=1}^L(x-p_i)$, so it is asymptotically sharp.  In the opposite
direction, the degree of an integer polynomial of height $H\ge2$ all of whose
roots are nonzero rationals is at most a constant times $\log H$~\cite{htv},
and the minimal height of a degree-$n$ polynomial over a fixed number field
with all roots in $K^\times$ grows exponentially in $n$~\cite{tromp}; the
bound here constrains neither degree nor height but the additive coefficient
mass.  The contribution
relative to reduced-height duality is the displaced-zero comparison and its
iteration; the normalized determinant sequence and its consecutive-zero
evaluation are from Dubickas and Jankauskas~\cite{dj}.  The confluent
Proposition~\ref{prop:confluent} is supplementary: it removes the
distinctness of the nodes from the certificate, but the distinct-node
theorem already suffices for the order-statistic and mass bounds.

\paragraph{Data availability.}
No datasets were generated or analyzed in this study.

\begin{thebibliography}{9}
\small
\bibitem{karlin-studden}
S.~Karlin and W.~J. Studden,
\newblock \emph{Tchebycheff Systems: With Applications in Analysis and Statistics},
\newblock Interscience, New York, 1966.

\bibitem{dubickas}
A.~Dubickas,
\newblock Arithmetical properties of powers of algebraic numbers,
\newblock \emph{Bulletin of the London Mathematical Society} 38 (2006), 70--80,
\newblock \href{https://doi.org/10.1112/S0024609305017728}
{doi:10.1112/S0024609305017728}.

\bibitem{schinzel}
A.~Schinzel,
\newblock On the reduced length of a polynomial with real coefficients,
\newblock \emph{Functiones et Approximatio} 35 (2006), 271--306,
\newblock \href{https://doi.org/10.7169/facm/1229442629}
{doi:10.7169/facm/1229442629}.

\bibitem{schinzel-complex}
A.~Schinzel,
\newblock The reduced length of a polynomial with complex coefficients,
\newblock \emph{Acta Arithmetica} 133 (2008), 73--81,
\newblock \href{https://doi.org/10.4064/AA133-1-5}
{doi:10.4064/AA133-1-5}.

\bibitem{bombieri-vaaler}
E.~Bombieri and J.~D. Vaaler,
\newblock Polynomials with low height and prescribed vanishing,
\newblock in A.~C. Adolphson, J.~B. Conrey, A.~Ghosh, and R.~I. Yager (eds.),
\emph{Analytic Number Theory and Diophantine Problems},
Progress in Mathematics 70, Birkh\"auser, 1987, 53--73,
\newblock \href{https://doi.org/10.1007/978-1-4612-4816-3_3}
{doi:10.1007/978-1-4612-4816-3\_3}.

\bibitem{dj}
A.~Dubickas and J.~Jankauskas,
\newblock On the reduced height of a polynomial,
\newblock \emph{Publicationes Mathematicae Debrecen} 71 (2007), 325--348,
\newblock \href{https://doi.org/10.5486/PMD.2007.3728}
{doi:10.5486/PMD.2007.3728}.

\bibitem{grt}
M.~Giesbrecht, D.~S. Roche, and H.~Tilak,
\newblock Computing sparse multiples of polynomials,
\newblock \emph{Algorithmica} 64 (2012), 454--480,
\newblock \href{https://doi.org/10.1007/s00453-012-9652-4}
{doi:10.1007/s00453-012-9652-4}.

\bibitem{bbbp-multiple}
F.~Beaucoup, P.~Borwein, D.~W. Boyd and C.~Pinner,
\newblock Multiple roots of $[-1,1]$ power series,
\newblock \emph{Journal of the London Mathematical Society} (2) 57 (1998), 135--147,
\newblock \href{https://doi.org/10.1112/S0024610798005857}
{doi:10.1112/S0024610798005857}.

\bibitem{bbbp-ray}
F.~Beaucoup, P.~Borwein, D.~W. Boyd and C.~Pinner,
\newblock Power series with restricted coefficients and a root on a given ray,
\newblock \emph{Mathematics of Computation} 67 (1998), 715--736,
\newblock \href{https://doi.org/10.1090/S0025-5718-98-00960-0}
{doi:10.1090/S0025-5718-98-00960-0}.

\bibitem{frw}
M.~Filaseta, M.~L. Robinson and F.~S. Wheeler,
\newblock The minimal Euclidean norm of an algebraic number is effectively
computable,
\newblock \emph{Journal of Algorithms} 16 (1994), 309--333,
\newblock \href{https://doi.org/10.1006/jagm.1994.1015}
{doi:10.1006/jagm.1994.1015}.

\bibitem{htv}
L.~Hajdu, R.~Tijdeman and N.~Varga,
\newblock On polynomials with only rational roots,
\newblock \emph{Mathematika} 69 (2023), 867--878,
\newblock \href{https://doi.org/10.1112/mtk.12209}
{doi:10.1112/mtk.12209}.

\bibitem{mignotte}
M.~Mignotte,
\newblock Sur les multiples des polyn\^omes irr\'eductibles,
\newblock \emph{Bulletin de la Soci\'et\'e Math\'ematique de Belgique} 27
(1975), 225--229.

\bibitem{polya-szego}
G.~P\'olya and G.~Szeg\H{o},
\newblock \emph{Problems and Theorems in Analysis II},
\newblock Springer, Berlin, 1976.

\bibitem{tromp}
T.~Tromp,
\newblock On the height of polynomials that split completely over a fixed
number field,
\newblock arXiv:2509.12064 (2025),
\newblock \href{https://doi.org/10.48550/arXiv.2509.12064}
{doi:10.48550/arXiv.2509.12064}.
\end{thebibliography}

\end{document}
